\section{From local isometries to an integral isometry}
\label{global:section}

We now turn the local and rational isometries of
Proposition~\ref{local:rational} into one integral isometry, after two
further operations.  A second tensor square removes all eigenspace
determinants from the discrepancy between the local and rational maps.
Two common polynomial summands then put the relevant isometry groups
within the scope of strong approximation.

Write the tensor-square data as
\[
 A_j=(M_j,\tau_j,b_j)
     =(L_j^{\otimes2},T_j^{\otimes2},B_j^{\otimes2}).
\]
Fix the equivariant rational isometry $R:A_1\longrightarrow A_2$ and
the equivariant local lattice isometries $D_\ell$ supplied by
Proposition~\ref{local:rational}, with $D_\ell=R$ for all but finitely
many primes $\ell$.

\subsection{Cancelling the determinant on each eigenspace}

The ratio $R^{-1}D_\ell$ preserves $b_1$ and commutes with $\tau_1$
over $\Q_\ell$, but need not have determinant one on each eigenspace.
To cancel these
determinants by a second tensor square, we use complex conjugation on
the real-coefficient polynomials $q_j$.  Recall that their common
ambient dimension $s$ is even and that $B_j=S_jT_j^r$.

\begin{lemma}[Integral reversers]\label{local:reverser}
There is an integral involution $C_j$ of $L_j$ such that
\[
 C_j^{-1}T_jC_j=T_j^{-1},\qquad C_j^{\mathsf t}B_jC_j=B_j.
\]
\end{lemma}

\begin{proof}
Complex conjugation preserves the argument-zero Milnor page and
exchanges its positive and negative normal directions; let $C_j$ be
its action on middle homology.  We check the sign directly in the
linking definition of the Seifert form~\cite{Milnor1968}.
The ambient sphere has dimension $2s-1$, and the cycles have dimension
$s-1$.  Complex conjugation has degree $(-1)^s=1$ on that sphere, and
interchanging the two linking cycles has sign $(-1)^{s^2}=1$.
If $y^+$ and $y^-$ denote the two push-offs, then
\[
 \begin{aligned}
 S_j(C_jx,C_jy)
 &=\operatorname{lk}(x,y^-)
  =\operatorname{lk}(y^-,x)\\
 &=\operatorname{lk}(y,x^+)=S_j(y,x).
 \end{aligned}
\]
The third equality moves both cycles along the page-normal flow.
Thus $C_j^{\mathsf t}S_jC_j=S_j^{\mathsf t}$.  Taking the monodromy
of this congruence gives $C_j^{-1}T_jC_j=T_j^{-1}$, and then
\[
 C_j^{\mathsf t}B_jC_j
   =S_j^{\mathsf t}T_j^{-r}=S_jT_j^{1-r}=B_j.
\]
The map on oriented cycles already includes conjugation's effect on
page orientation; the calculation uses the orientation of the ambient
sphere for linking.
\end{proof}

Tensoring gives an integral isometry
$C_j^{(2)}=C_j\otimes C_j$ of $b_j$ that reverses $\tau_j$.
On the fourth tensor powers of the original data define
\begin{equation}
\label{global:corrected-maps}
 \widehat R
 =R\otimes\bigl(C_2^{(2)}R(C_1^{(2)})^{-1}\bigr),
 \qquad
 \widehat D_\ell
 =D_\ell\otimes\bigl(C_2^{(2)}D_\ell(C_1^{(2)})^{-1}\bigr).
\end{equation}
Both second factors intertwine the operators: each conjugating map
reverses its operator, and $R$ and $D_\ell$ intertwine them.
They also preserve the forms.  Consequently $\widehat R$ is a rational
equivariant isometry, and $\widehat D_\ell$ is an equivariant isometry
of the fourth-power lattices over $\Z_\ell$.

\begin{lemma}[Eigenspace determinant cancellation]
\label{global:determinants}
For every prime $\ell$, the ratio
$\widehat R^{-1}\widehat D_\ell$ has determinant $1$ on every
eigenspace of $\tau_1\otimes\tau_1$ over $\overline{\Q_\ell}$.
This remains true after tensoring both maps with identity maps on any
common finite-order operator data.
\end{lemma}

\begin{proof}
Put $a=R^{-1}D_\ell$ and $C=C_1^{(2)}$.  Multiplying the two maps in
\eqref{global:corrected-maps} gives the exact identity
\begin{equation}
\label{global:ratio}
 \widehat R^{-1}\widehat D_\ell=a\otimes CaC^{-1}.
\end{equation}
In particular, no compatibility of $R$ with the reversers is required.

Let $W_\gamma$ be the $\gamma$-eigenspace of $\tau_1$ over
$\overline{\Q_\ell}$, and set
\[
 m_\gamma=\dim W_\gamma,
 \qquad t_\gamma=\det(a|_{W_\gamma}).
\]
Invariance and nondegeneracy of $b_1$ give a perfect pairing between
$W_\gamma$ and $W_{\gamma^{-1}}$.  Since $a$ is an isometry,
\begin{equation}
\label{global:inverse-determinants}
 m_{\gamma^{-1}}=m_\gamma,
 \qquad t_{\gamma^{-1}}=t_\gamma^{-1}.
\end{equation}
For $\gamma=1$, this says $t_1^2=1$; no choice of its sign is needed.
The reverser carries $W_{\eta^{-1}}$ onto $W_\eta$, so the determinant
of $CaC^{-1}$ on $W_\eta$ is $t_{\eta^{-1}}$.

For each $\lambda$, its entire eigenspace in the tensor square is
\[
 \bigoplus_{\gamma\eta=\lambda}W_\gamma\otimes W_\eta,
\]
where the sum ranges over all ordered pairs of occurring eigenvalues.
By the determinant formula for a tensor product, the determinant of
\eqref{global:ratio} on this space is
\begin{align*}
 \prod_{\gamma\eta=\lambda}
       t_\gamma^{m_\eta}t_{\eta^{-1}}^{m_\gamma}
 &=
 \frac{\displaystyle\prod_{\gamma\eta=\lambda}t_\gamma^{m_\eta}}
      {\displaystyle\prod_{\gamma\eta=\lambda}t_\eta^{m_\gamma}}
 =1.
\end{align*}
The last equality exchanges $\gamma$ and $\eta$ in the denominator.
It includes diagonal pairs and all coincidences among product
eigenvalues.  Finally, tensoring with an identity gives, on each new
eigenspace, a product of powers of the already unit determinants.
\end{proof}

\subsection{Two common summands}

Use new disjoint variables to form
\begin{equation}
\label{global:common-polynomials}
 u=z_1^3+\cdots+z_6^3,
 \qquad v=z_7^{11}+z_8^{13}.
\end{equation}
Both have an even number of variables and odd-order monodromy.  Denote
their symmetric-form and operator data by $\mathcal U$ and
$\mathcal V$, and put
\begin{equation}
\label{global:final-data}
 H_j=(\Lambda_j,\Theta_j,\beta_j)
     =A_j^{\otimes2}\otimes\mathcal U\otimes\mathcal V.
\end{equation}
Extend the maps in \eqref{global:corrected-maps} by the identity on
these two common factors, keeping the notation $\widehat R$ and
$\widehat D_\ell$.

The six cubes will make each eigenspace form indefinite.  The degrees
$11$ and $13$ exclude real eigenvalues from the final operator, even
though the preceding tensor squares may have acquired eigenvalue $1$.
These two primes also enter the later proof that the links are homotopy
spheres in Lemma~\ref{topo:sphere}.

\begin{lemma}[Spectral properties of the added factors]
\label{global:indefinite}
Every eigenvalue of $\Theta_j$ has odd order divisible by both $11$
and $13$.  On every complex eigenspace, the hermitian form
\[
 (x,y)\longmapsto\beta_j(x,\overline y)
\]
is nondegenerate and indefinite; the complex dimension of that
eigenspace is at least $21$.
\end{lemma}

\begin{proof}
Before adjoining $\mathcal V$, the operator order divides
$3^a\cdot5\cdot7$ for some $a$.  The eigenvalues of $v$ are
\[
 \zeta_{11}^b\zeta_{13}^c,
 \qquad 1\le b\le10,\quad 1\le c\le12.
\]
Their nontrivial $11$- and $13$-primary components cannot cancel
against an eigenvalue of the earlier tensor factors.  Thus every
final order contains both primes and remains odd.  In particular,
neither $1$ nor $-1$ is a final eigenvalue.

We calculate the signatures supplied by $u$.  Let
\[
 A=\begin{pmatrix}1&-1\\0&1\end{pmatrix}.
\]
In the convention of Theorem~\ref{bg:ts}, a one-cube Seifert matrix
is $S_c=-A$~\cite[Corollary~3]{Sakamoto1974}.  Its chosen monodromy is
\[
 t=-A^{-1}A^{\mathsf t}
   =\begin{pmatrix}0&-1\\1&-1\end{pmatrix}.
\]
Here the minus sign is the one-variable convention.  The form
\[
 J=At^2=\begin{pmatrix}0&1\\-1&0\end{pmatrix}
\]
is alternating and $t$-invariant.  Iterating Theorem~\ref{bg:ts}
gives $S_u=(-1)^{15}(-A)^{\otimes6}=-A^{\otimes6}$.
The sixfold monodromy is $t^{\otimes6}$, and
$2\cdot2\equiv1\pmod3$, so its symmetric form is
\[
 B_u=S_u(t^{\otimes6})^2=-J^{\otimes6}.
\]
We compute the signatures of $J^{\otimes6}$ below; passing to $B_u$
reverses every sign.

Set $\omega=e^{2\pi i/3}$ and $e_+=(1,-\omega)$,
$e_-=(1,-\omega^2)$.  These are eigenvectors of $t$ with
eigenvalues $\omega$ and $\omega^2$, respectively, and
\[
 J(e_+,\overline{e_+})=i\sqrt3,
 \qquad J(e_-,\overline{e_-})=-i\sqrt3,
 \qquad J(e_+,\overline{e_-})=0.
\]
In the resulting orthogonal tensor basis, a vector with $k$ copies
of $e_+$ has eigenvalue $\omega^{-k}$ and length
$27(-1)^{k+1}$; there are $\binom6k$ such vectors.  We obtain the
following signatures, up to reversal of every sign:
\[
\begin{array}{c|c|c|c}
 \text{eigenvalue}&k&\text{signature}&\text{dimension}\\\hline
 1&0,3,6&(20,2)&22\\
 \omega&2,5&(6,15)&21\\
 \omega^2&1,4&(6,15)&21
\end{array}
\]
Every eigenspace of this common factor is therefore indefinite.

For any real nondegenerate symmetric form $b$ invariant under a
finite-order operator, its complex bilinear extension pairs the
$\gamma$-eigenspace perfectly with the $\gamma^{-1}$-eigenspace.
Consequently $b(x,\overline y)$ restricts to a nondegenerate
hermitian form on each eigenspace.  Distinct eigenspaces are
orthogonal for this hermitian form.

Apply this observation to each of the three factors in
\eqref{global:final-data}.  A final eigenspace is the orthogonal
direct sum of the tensor products of factor eigenspaces whose
eigenvalues have the prescribed product.  Each nonzero summand
contains one of the indefinite spaces in the table.  Diagonalizing
the other two hermitian forms shows that its tensor form has both
signs: multiply any nonzero diagonal entry from the other factors
by a positive and a negative entry from the six-cube factor.
Every summand is thus indefinite and has dimension at least $21$.
The same conclusions hold for their orthogonal direct sum.
\end{proof}

\subsection{The special unitary group and strong approximation}

\begin{proposition}[Integral equivariant isometry]
\label{global:integral}
The triples $H_1$ and $H_2$ in \eqref{global:final-data} are
isometric over $\Z$: there is a lattice isomorphism
$U:\Lambda_1\longrightarrow\Lambda_2$ satisfying
\[
 U\Theta_1=\Theta_2U,
 \qquad U^{\mathsf t}\beta_2U=\beta_1.
\]
\end{proposition}

\begin{proof}
Put $V=\Lambda_1\otimes\Q$.  Since $\Theta_1$ has finite order,
its minimal polynomial is a product of distinct cyclotomic
polynomials.  Decompose the rational operator space as
\[
 V=\bigoplus_{d\in\mathcal D}V_d,
 \qquad
 E_d=\Q(\zeta_d),\quad F_d=\Q(\zeta_d+\zeta_d^{-1}),
 \quad n_d=\dim_{E_d}V_d.
\]
The action of $\Theta_1$ gives $V_d$ its $E_d$-vector space
structure.  Distinct $V_d$ are orthogonal for $\beta_1$.
By Lemma~\ref{global:indefinite}, all $d$ are odd and divisible
by $11\cdot13$, so $E_d/F_d$ is quadratic with nontrivial
conjugation.  The adjoint of multiplication by $a\in E_d$ is
multiplication by $\overline a$, because the adjoint of $\Theta_1$
is $\Theta_1^{-1}$.

For completeness, the nondegenerate trace pairing of $E_d/\Q$
defines a unique $E_d$-valued form $h_d$ by
\[
 \tr_{E_d/\Q}\bigl(a h_d(x,y)\bigr)=\beta_1(ax,y)
 \quad(a\in E_d).
\]
This form is linear in its first argument, conjugate-linear in
its second, and satisfies $h_d(y,x)=\overline{h_d(x,y)}$.
These identities follow by testing against every $a$ in the
displayed trace equation.  The same equation proves its
nondegeneracy.  It also shows that an $E_d$-linear map preserves
$\beta_1$ exactly when it preserves $h_d$.

Let $G$ be the algebraic $\Q$-group of equivariant self-isometries of
$(V,\Theta_1,\beta_1)$ with determinant $1$ on every eigenspace
over an algebraic closure.  These conditions are defined over
$\Q$.  The preceding description identifies this algebraic group as
\begin{equation}
\label{global:unitary-group}
 G=\prod_{d\in\mathcal D}\Res_{F_d/\Q}\SU(h_d).
\end{equation}
Indeed, commuting with $\Theta_1$ is exactly $E_d$-linearity on
each $V_d$.  Each embedding of $E_d$ into an algebraic closure
gives one $\Theta_1$-eigenspace, and the determinant on it is
the corresponding image of the $E_d$-determinant.  Requiring
all these determinants to be $1$ gives precisely $\SU(h_d)$.
In particular this requirement removes the unitary determinant
tori, and the absence of eigenvalues $\pm1$ leaves no
orthogonal factor.

Every embedding $F_d\hookrightarrow\R$ extends to a conjugate
pair of embeddings of $E_d$ into $\C$.  At this real place the trace
identity writes the underlying real form as $2\operatorname{Re}h_d$.
After complexification, its pairing $(x,y)\mapsto\beta_1(x,\overline y)$
on either corresponding operator eigenspace is the associated complex
hermitian form.  Lemma~\ref{global:indefinite} therefore makes it
indefinite, and $n_d\ge21$.
The group $\SU(h_d)$ is connected, simply connected, and
absolutely almost simple over $F_d$, since over an algebraic
closure it becomes $\mathrm{SL}_{n_d}$
\cite[Theorem~24.44, Remark~24.46, and Summary~24.66]{Milne2022}.  At each real place it
is a noncompact group $\SU(p,q)$ with $p,q>0$.

Theorem~\ref{bg:approximation}, applied to $\SU(h_d)$ over $F_d$
with all archimedean places omitted, now gives density of
$\SU(h_d)(F_d)$ in its finite adelic group
\cite[Theorem~A]{Platonov1969}, with~\cite{Platonov1970}; see also
\cite[Theorem~2.3]{Rapinchuk2014}.
Restriction of scalars and the finite product in
\eqref{global:unitary-group} therefore give
\begin{equation}
\label{global:density}
 \overline{G(\Q)}=G(\A_{\Q,f}).
\end{equation}
Here finite adeles mean the restricted product over all finite
places.  No isotropy over the number field itself is being
assumed; the required noncompactness is supplied at the real
places just checked.

We spell out how \eqref{global:density} gives an integral map.
For each prime set
\[
 a_\ell=\widehat R^{-1}\widehat D_\ell,
 \qquad \Lambda_{j,\ell}=\Lambda_j\otimes\Z_\ell.
\]
Lemma~\ref{global:determinants} gives $a_\ell\in G(\Q_\ell)$,
and $a_\ell=1$ for all but finitely many $\ell$.  Thus
$a=(a_\ell)$ is a finite adelic point.  Define
\[
 K_\ell=\{g\in G(\Q_\ell):g\Lambda_{1,\ell}
                   =\Lambda_{1,\ell}\},
 \qquad K=\prod_\ell K_\ell.
\]
In a basis of the actual lattice $\Lambda_1$, the subgroup
$K_\ell$ is $G(\Q_\ell)\cap\GL_{\operatorname{rank}\Lambda_1}(\Z_\ell)$.
It is compact and open in $G(\Q_\ell)$. To identify their product with
an adelic open subgroup, take the schematic closure of $G$ in
$\GL(\Lambda_1)$ over $\Z$. After inverting finitely many primes this
is an integral model of the same affine group, and its
$\Z_\ell$-points are exactly $K_\ell$. Thus the $K_\ell$ agree with
the reference integral subgroups defining the restricted product at
all but finitely many primes. Their product $K$ is consequently compact
and open in $G(\A_{\Q,f})$.  This definition uses the full lattice and
requires no integral splitting of its cyclotomic components.

By \eqref{global:density}, the nonempty open coset $aK$ contains
a point $g\in G(\Q)$.  For every prime we can write
$g=a_\ell k_\ell$ with $k_\ell\in K_\ell$, and hence
\[
 (\widehat Rg)(\Lambda_{1,\ell})
 =\widehat R a_\ell(\Lambda_{1,\ell})
 =\widehat D_\ell(\Lambda_{1,\ell})
 =\Lambda_{2,\ell}.
\]
A full integral lattice satisfies
\[
 \Lambda_j
 =\{x\in\Lambda_j\otimes\Q:
                x\in\Lambda_{j,\ell}\text{ for every prime }\ell\}:
\]
in a lattice basis this is the elementary identity
$\Q\cap\bigcap_\ell\Z_\ell=\Z$.
Applying this characterization to $\widehat Rg$ and its inverse
shows that $U=\widehat Rg$ carries $\Lambda_1$ onto $\Lambda_2$.
It is an equivariant isometry because both factors are such
isometries, which proves the proposition.
\end{proof}

\begin{remark}
\label{global:recover-seifert}
An equivariant isometry of the symmetric data recovers an isometry
of the original Seifert data.  Choose one integer $r$ satisfying
$2r\equiv1$ modulo a common odd exponent of the final operators.
If $S_j'=\beta_j\Theta_j^{-r}$, then the isomorphism in
Proposition~\ref{global:integral} satisfies
\[
 U^{\mathsf t}S_2'U
 =U^{\mathsf t}\beta_2U\Theta_1^{-r}
 =\beta_1\Theta_1^{-r}=S_1'.
\]
The tensor-product convention can multiply both final forms by
the same sign, which does not alter this identity.
\end{remark}
